\chapter{Quatrex Documentation}
\label{ch:quatrex_doc}

% ----------------------------------------------------------------------
\subsection{QuaTrEx Input format}
\label{sec:quatrex_input}
% ----------------------------------------------------------------------

The QuaTrEx solver reads three input files which
specify the real-space Hamiltonian, the crystal structure, and the
transport parameters.

\subsubsection{Input files}
\label{sec:input_files}

\paragraph{\texttt{dynamical\_matrix.mat}}
This MATLAB \texttt{.mat} file stores the real-space blocks
$H(\vc{R})$ in Convention~B, already converted to
$(\mathrm{rad/s})^2$.  Each block is stored under a string key of the
form
\begin{equation}
  \texttt{"[nx, ny, nz]"},
\end{equation}
mapping to a real matrix of dimension
$(n_\mathrm{dof}\times n_\mathrm{dof})$, where
$n_\mathrm{dof}=3N_b$ and $N_b$ is the number of atoms in the unit
cell.  For nearest-neighbour interactions extracted on a
$3\times3\times3$ mesh, the file contains up to 27 blocks.

The solver interprets these blocks as real-space couplings and, for
periodic transverse directions, performs the Fourier sum
\begin{equation}
  D(\vc{k}_\perp)
  = \sum_{(n_x,n_y)} M(n_x,n_y)\;
    e^{2\pi i\,(k_x n_x + k_y n_y)}.
\end{equation}

\paragraph{\texttt{structure.xyz}}
This is an extended XYZ file containing the lattice vectors in the
header and Cartesian atomic positions in the body.  It follows the ASE
extended-XYZ convention, for example
\begin{verbatim}
8
Lattice="5.4018 0.0 0.0 0.0 5.4018 0.0 0.0 0.0 5.4018" ...
Si  0.00000000  0.00000000  0.00000000
...
\end{verbatim}
The loader (\texttt{distributed\_read\_xyz}) extracts the $3\times3$
lattice matrix and Cartesian atom positions.  The number of atoms times
\texttt{num\_orbitals\_per\_atom} determines $n_\mathrm{dof}$.

\paragraph{\texttt{quatrex\_config.toml}}
This file specifies the transport setup, such as transport direction,
number of transport cells, neighbour cutoff, transverse $k$-mesh,
broadening, and the open-boundary-condition algorithm.  Important fields are
\begin{center}
\begin{tabular}{lll}
  \toprule
  Field & Example value & Purpose \\
  \midrule
  \texttt{transport\_direction}       & \texttt{'z'}            & transport axis \\
  \texttt{construct\_from\_unit\_cell}& \texttt{true}           & input mode selector \\
  \texttt{num\_transport\_cells}      & 4                       & device length in unit cells \\
  \texttt{neighbor\_cell\_cutoff}     & \texttt{[1,1,1]}        & block-range cutoff \\
  \texttt{kpoint\_grid}               & \texttt{[4,4,1]}        & transverse mesh \\
  \texttt{kpoint\_shift}              & \texttt{[0,0,0]}        & Monkhorst--Pack shift \\
  \texttt{num\_orbitals\_per\_atom.Si}& 3                       & DOFs per atom \\
  \texttt{phonon.eta}                 & \texttt{1e-8}           & broadening $(\mathrm{rad/s})^2$ \\
  \texttt{phonon.obc.algorithm}       & \texttt{"sancho-rubio"} & OBC method \\
  \bottomrule
\end{tabular}
\end{center}

\subsubsection{Input modes}
\label{sec:input_modes}

QuaTrEx supports two input modes, selected by
\texttt{construct\_from\_unit\_cell}.  In both modes the input is a
MATLAB \texttt{.mat} file loaded by \texttt{scipy.io.loadmat}.  The
loader (\\\texttt{distributed\_load} in
\texttt{qttools/utils/mpi\_utils.py}) parses keys that start with
\texttt{[} and interprets them as integer tuples:
\begin{lstlisting}
arr = loadmat(path)
arr = {
    tuple(map(int, r.strip("[]").split(","))): h_r
    for r, h_r in arr.items()
    if r.startswith("[")
}
\end{lstlisting}
That is, a key string such as \texttt{"[0, 1, -1]"} becomes the Python
tuple \texttt{(0, 1, -1)}.  Each value is a 2D NumPy array (real or
complex) representing a matrix block.

\paragraph{Mode 1: pre-assembled device matrix}
If \texttt{construct\_from\_unit\_cell = false}, the
\texttt{dynamical\_matrix.mat} file contains a single pre-assembled
device matrix stored under the key \texttt{(0,0,0)}:
\begin{lstlisting}
matrix_sparray = distributed_load(f"{matrix_name}.mat")
assert (0, 0, 0) in matrix_sparray.keys()
matrix_sparray = matrix_sparray[(0, 0, 0)]
\end{lstlisting}
The result is a single sparse matrix of size
$(N_\mathrm{blocks}\times B)^2$, where $B$ is the block size and
$N_\mathrm{blocks}$ is the number of transport cells.  No transverse
$k$-point Fourier assembly is performed.

\paragraph{Mode 2: unit-cell blocks}
If \texttt{construct\_from\_unit\_cell = true}, the file contains the
dictionary of real-space unit-cell blocks
$H(n_x,n_y,n_z)$.  The solver loads all blocks, constructs a dense
translation grid, optionally trims it using
\texttt{neighbor\_cell\_cutoff}, expands the blocks into
transport-direction BTD matrices, and performs the transverse
Fourier assembly on the chosen Monkhorst--Pack mesh.

The function \texttt{\_load\_matrix\_from\_unit\_cell} loads all keys,
determines the bounding box from the minimum and maximum key
coordinates, and packs the blocks into a 5D array:
\begin{lstlisting}
# keys: array of all (nx, ny, nz) tuples
min_coords = keys.min(axis=0)
max_coords = keys.max(axis=0)
grid_shape = max_coords - min_coords + 1
unit_cells = zeros(grid_shape + matrix_shape)
for coord, matrix in matrices.items():
    unit_cells[coord] = matrix
\end{lstlisting}
The resulting array has shape
$(G_x,G_y,G_z,n_\mathrm{dof},n_\mathrm{dof})$, where
$(G_x,G_y,G_z)$ is the grid shape (e.g.\ $(3,3,3)$ for
$\{-1,0,+1\}^3$).  The loader checks that
$G_x \times G_y \times G_z = \text{number of keys}$, i.e.\ the grid
must be dense: no missing entries are permitted.  Missing blocks must
be stored as zero matrices.

If \texttt{neighbor\_cell\_cutoff} is set (e.g.\ \texttt{[1,1,1]}),
\texttt{trim\_tight\_binding\_matrix} zeros out all blocks beyond that
range.  For cutoff $(c_x,c_y,c_z)$, a block at index
$(n_x,n_y,n_z)$ (centred at zero) is kept only if
$|n_x| \le c_x$, $|n_y| \le c_y$, $|n_z| \le c_z$.

The function \texttt{expand\_tight\_binding\_matrix} converts the
unit-cell blocks into a BTD sparse matrix for a given transverse
periodic shift.  Along the transport direction $z$:
\begin{enumerate}
  \item The supercell size in the transport direction is $G_z / 2$
        (integer division), and 1 in the transverse directions.  For
        $G_z = 3$, this gives a supercell of one unit cell
        ($3 // 2 = 1$).
  \item Three BTD blocks are extracted for each transverse periodic
        shift $(n_x,n_y)$:
        \begin{itemize}
          \item $n_z=-1$: lower off-diagonal,
          \item $n_z=0$: diagonal,
          \item $n_z=+1$: upper off-diagonal.
        \end{itemize}
  \item Each block is assembled by \texttt{\_get\_transport\_block},
        which loops over all pairs of local shifts within the supercell
        and selects the appropriate unit-cell block for each pair.
  \item The three blocks are tiled $N_\text{transport}$ times along the
        diagonal to build the full BTD sparse matrix.
\end{enumerate}

The function \texttt{\_assemble\_kpoint} then performs the Fourier sum
over transverse periodic shifts.  For each $\vc{k}_\perp$ on a
Monkhorst--Pack mesh:
\begin{equation}\label{eq:assemble}
  D(\vc{k}_\perp)
  = \sum_{(n_x,n_y)} M(n_x,n_y)\;
    e^{2\pi i\,(k_x n_x + k_y n_y)},
\end{equation}
where $M(n_x,n_y)$ is the BTD sparse matrix for periodic shift
$(n_x,n_y)$, already expanded in the transport direction.  The
$k$-points on the Monkhorst--Pack mesh are
\begin{equation}
  k_i = \frac{n_i + 0.5}{N_i} - 0.5 + \text{shift}_i,
  \qquad n_i = 0,\ldots,N_i-1.
\end{equation}
The result is a stack of BTD matrices indexed by
$(\omega,k_x,k_y)$.

% ----------------------------------------------------------------------
\subsection{Heterogeneous-device extension}
\label{sec:heterogeneous}
% ----------------------------------------------------------------------

The original harmonic unit-cell input mode in QuaTrEx constructs a
homogeneous device by repeating the same block set across all transport
cells.  To model interfaces such as Si/Ge, the harmonic assembly was
extended to support piecewise-defined left, device, and right regions
with distinct block dictionaries.

\subsubsection{L|D|R device model}

The transport region is partitioned into three sections,
\begin{equation}\label{eq:ldr}
  \underbrace{L\;\cdots\;L}_{n_L\;\text{cells}}
  \;\underbrace{D\;\cdots\;D}_{n_D\;\text{cells}}
  \;\underbrace{R\;\cdots\;R}_{n_R\;\text{cells}},
  \qquad
  n_L + n_D + n_R = n_\mathrm{transport},
\end{equation}
where the left and right parts serve as lead-like material regions and
the central part represents the interface or device region.

\subsubsection{Configuration changes}

To support this heterogeneous harmonic setup, four new configuration
fields were added to \texttt{DeviceConfig}:
\begin{center}
\begin{tabular}{lll}
  \toprule
  Field & Type & Purpose \\
  \midrule
  \texttt{num\_left\_cells}  & \texttt{PositiveInt | None} & number of left cells $n_L$ \\
  \texttt{num\_right\_cells} & \texttt{PositiveInt | None} & number of right cells $n_R$ \\
  \texttt{left\_matrix}      & \texttt{str | None}         & left-lead \texttt{.mat} file \\
  \texttt{right\_matrix}     & \texttt{str | None}         & right-lead \texttt{.mat} file \\
  \bottomrule
\end{tabular}
\end{center}
This mode requires
\texttt{construct\_from\_unit\_cell = true} and must satisfy
\begin{equation}
  n_L + n_R < n_\mathrm{transport},
\end{equation}
so that at least one transport cell remains in the device region.

\subsubsection{Matrix assembly}

The heterogeneous assembly is implemented by extending the harmonic
unit-cell expansion to allow different block dictionaries in different
transport cells.  Two new functions were added to
\texttt{quatrex/device/inputs.py}.

\texttt{\_expand\_heterogeneous\_matrix} builds a BTD sparse matrix for
a single transverse periodic shift $(n_x,n_y)$.  Each cell is assigned
to one of the three regions according to its index.  The diagonal block
of cell $i$ is selected according to
\begin{equation}
  D_{ii} =
  \begin{cases}
    D_L, & i < n_L, \\
    D_D, & n_L \le i < n_L+n_D, \\
    D_R, & i \ge n_L+n_D.
  \end{cases}
\end{equation}

The routine \func{_create_heterogeneous_matrix} constructs the final
transport object. It loops over all transverse periodic shifts, calls
\func{_expand_heterogeneous_matrix} for each shift, and combines the
resulting blocks. At $\Gamma$, the transverse contributions are also
summed to determine the sparsity pattern.

\subsubsection{Boundary coupling}

At the interfaces between the $L$, $D$, and $R$ regions, the
off-diagonal couplings must be chosen consistently with Hermiticity.
For harmonic real-space blocks, the relevant relation is
\begin{equation}\label{eq:hermiticity_perp}
  H(\vc{R}_\perp,n_z)^\dagger
  = H(-\vc{R}_\perp,-n_z).
\end{equation}
At the left/device boundary, the super-diagonal coupling is taken from
the left-lead block $H_{01}^L$, while the corresponding sub-diagonal is
$H_{10}^L$.  At the device/right boundary, the couplings are taken from
the right-lead blocks $H_{01}^R$ and $H_{10}^R$.  This ensures that the
assembled heterogeneous matrix remains consistent with the physical lead
couplings on either side of the interface.
